\documentclass{article}
\usepackage{microtype,graphicx,booktabs,amsmath,amssymb}
\usepackage{hyperref}
\usepackage{mlsys2025}
% Content-only notice override for this internal draft; vendor style untouched.
\makeatletter
\renewcommand{\Notice@String}{Development draft for author review. Not submitted to MLSys.}
\makeatother
\hypersetup{hidelinks,pdfsubject={PageGauge MLSys 2027 development evidence}}
\mlsystitlerunning{PageGauge: Development Evidence}
\begin{document}
\twocolumn[
\mlsystitle{PageGauge: Factoring Affine KV Quantization\texorpdfstring{\\}{: }Development Evidence for an MLSys 2027 Draft}
\begin{mlsysauthorlist}
\mlsysauthor{Anonymous Authors}{anon}
\end{mlsysauthorlist}
\mlsysaffiliation{anon}{Anonymous Institution}
\mlsyscorrespondingauthor{Anonymous Author}{anon@example.org}
\vskip 0.3in
\begin{abstract}
This development report separates PageGauge's algebraic identity from
quantization fidelity and runtime efficiency. It records completed experiments
on an RTX 5090, including native low-bit cache comparisons and a matched
scale-placement control. The current evidence supports a narrow systems claim,
not universal quality or speed superiority. All quality results below use
exposed development windows; independent final evaluation, broader generalization,
and serving experiments remain outstanding. This is an evidence draft for
author review, not a submission-ready paper.
\end{abstract}
]
\printAffiliationsAndNotice{}

\section{Mathematical Contract}
For one query $q$ and KV head, let page $p$ contain reconstructed affine keys
and values
\begin{align}
 \widehat K_p &= \mathbf 1 c_K^\top+s_p^K Z_p^K,\\
 \widehat V_p &= \mathbf 1 c_V^\top+s_p^V Z_p^V.
\end{align}
The centers $c_K,c_V$ are common across the merged domain; each scale is a
scalar for its page and head. With $\alpha=1/\sqrt d$, the term
$\alpha q^\top c_K$ shifts every logit equally and cancels in softmax.
Let $\mathcal Q$ be quantized pages, $\mathcal E$ exact rows, and $P$ globally
normalized attention. With row-vector weights $P_p$ for each quantized page,
\begin{equation}
\begin{split}
 o^\top={}&c_V^\top+\sum_{p\in\mathcal Q}s_p^V P_p Z_p^V\\
          &+\sum_{j\in\mathcal E}P_j(v_j-c_V)^\top.
\end{split}
\end{equation}
Exact regions use centered keys and values and join the same online softmax
merge. The value center is restored once after the merge, not once per page.
This identity is exact in real arithmetic relative to the reconstructed mixed
cache, not the original unquantized model. Finite-precision execution and
quantization error therefore require distinct measurements.

\section{Quality and Storage}
Table~\ref{tab:quality} uses Mistral-7B-v0.3, batch one, 20,480 initial tokens,
and 1,536 recurrent teacher-forced steps. Eight disjoint exposed TRAIN windows
provide 12,288 scored labels. Perplexity (PPL) is computed by aggregating token
negative log likelihood before exponentiation. Each row retains its own HF
reference: PageGauge/FlashInfer use Transformers 4.57.6 and SDPA; the native
KIVI and BitDecoding integration use Transformers 4.36.2 and eager attention.
Actual token IDs are matched across stacks.

\begin{table*}[t]
\caption{Development quality and served KV storage. PPL ratios compare each
implementation with its own HF reference. Intervals are descriptive paired-window
bootstrap intervals, not independent-book or prospective non-inferiority tests.
Cache storage includes metadata and exact residual regions but excludes model
weights, activations, allocator reserve, and separate staging.}
\label{tab:quality}
\centering\small
% Generated development evidence; do not edit by hand.
\begin{tabular}{lrrrrrr}
\toprule
Method & PPL & HF PPL & PPL/HF & Descriptive 95\% interval & KV MiB & KV reduction\\\midrule
FlashInfer FP16 & 5.20421 & 5.20410 & 1.00002 & [0.99997, 1.00008] & 2752.00 & 0.00\% \\
PageGauge INT8 & 5.20466 & 5.20410 & 1.00011 & [0.99993, 1.00027] & 1737.72 & 36.86\% \\
KIVI INT4 & 5.20767 & 5.20417 & 1.00067 & [0.99975, 1.00186] & 861.38 & 68.70\% \\
KIVI INT2 & 5.35293 & 5.20417 & 1.02858 & [1.02312, 1.03493] & 517.62 & 81.19\% \\
BitDecoding INT4$^*$ & 5.20600 & 5.20417 & 1.00035 & [0.99984, 1.00115] & 876.00 & 68.17\% \\
NSN unquantized control & 5.20367 & 5.20426 & 0.99988 & [0.99978, 0.99997] & 2752.00 & 0.00\% \\
NSNQuant INT2 & 5.25094 & 5.20426 & 1.00897 & [1.00724, 1.01109] & 385.17 & 86.00\% \\
\bottomrule
\end{tabular}

\par\smallskip\footnotesize
$^*$Native SM120-port kernels integrated into a shared Mistral harness,
not the upstream BitDecoding model engine. An additional 5.00012 MiB of
persistent staging is excluded from served KV storage and must be counted in
deployment capacity. NSNQuant includes resident quantizer/codebook buffers;
its unquantized control applies the native value/output rotation. NSN uses
Transformers 4.48.1 with full-prefix SDPA prefill and its own HF reference.
MiB denotes $2^{20}$ bytes.
\end{table*}

\begin{table*}[t]
\caption{Instrumented GPU self-time in ms per decode step, batch one, final
129 steps of the same 1,536-step development trajectory. The FP16 attention
column is the full cache for FlashInfer but only exact regions for PageGauge.
The total also includes operations not separately shown.}
\label{tab:profile}\centering\small
% Generated development evidence; do not edit by hand.
\begin{tabular}{lrrrrr}
\toprule
Method & Projections/LM & INT8 history & FP16 attention & Merge & Total\\\midrule
FlashInfer & 9.179 & 0.000 & 2.455 & 0.051 & 12.474 \\
PG split128 & 9.202 & 1.535 & 0.958 & 0.139 & 12.660 \\
PG split32 & 9.201 & 1.536 & 0.446 & 0.139 & 12.148 \\
\bottomrule
\end{tabular}

\end{table*}

The measured PageGauge PPL is close to its FP16 reference on this cohort, but
the four-bit competitors also have intervals containing one while using much
less cache storage. These observations do not establish PageGauge's predictive
quality superiority. Top-1 agreement and logit cosine remain supporting
diagnostics; neither replaces perplexity or generated-task evaluation.

\section{What Does Factorization Save?}
The scale-placement control uses identical INT8 history, exact regions, and
common-center factoring in both arms. Only placement of affine scales relative
to the FP16 code contractions changes. Eight fresh processes in ABBA/BAAB order
cover two TRAIN fixtures at batch four, initial context 20,480, and decode
length 1,536. Execution checks cover recurrence, graph replay, cache finalization,
new-history consumption, and sampled device exclusivity.

\begin{table}[t]
\caption{Scale-placement control: register-control latency divided by
factorized latency. Hierarchical intervals resample fixture then adjacent pair;
only two fixture clusters limit generality.}
\label{tab:causal}\centering\small
% Generated development evidence; do not edit by hand.
\begin{tabular}{lrr}
\toprule
Mode & Ratio & 95\% interval\\\midrule
Cache-neutral & 1.00433 & [1.00352, 1.00515] \\
Cache-hot & 1.00372 & [1.00289, 1.00456] \\
\bottomrule
\end{tabular}

\end{table}

The cache-neutral gain is approximately 0.43\%, not the entire previously
measured FlashInfer-relative advantage. The shared-center contribution has not
been isolated by this contrast. A stronger causal account must retain this
distinction rather than attributing all cache-format and scheduling effects to
scale placement.

\section{Exact-Region Scheduling}
An instrumented batch-one trajectory identifies a separate runtime cost:
the exact FP16 regions and merge operations can offset savings in compressed
history attention. Changing only the exact-region split from 128 to 32 pages
reduces its profiled cost (Table~\ref{tab:profile}); history work and projection
work remain similar. The number of launches does not decrease. Profiled kernel
self-time is not end-to-end latency. A separate optimized eight-process
ABBA/BAAB confirmation at batch four measures approximately 2.4\% lower model-step
latency for split32 (Table~\ref{tab:split}), with identical served cache storage
of 7,288,520,704 bytes. All eight execution assessments pass. This is a
PageGauge-versus-PageGauge scheduling contrast on two TRAIN fixtures, not a new
FlashInfer-relative or independent-quality claim.

\begin{table}[t]
\caption{Optimized exact-region scheduling: split128 latency divided by
split32 latency at B4/C20480/D1536. Eight fresh processes, four adjacent pairs,
two fixture clusters; hierarchical 95\% intervals.}
\label{tab:split}\centering\small
% Generated development evidence; do not edit by hand.
\begin{tabular}{lrr}
\toprule
Mode & Ratio & 95\% interval\\\midrule
Cache-neutral & 1.02464 & [1.02406, 1.02535] \\
Cache-hot & 1.02486 & [1.02441, 1.02520] \\
\bottomrule
\end{tabular}

\end{table}

\section{Second Model Family: Qwen3-8B}
Eight exposed WikiText TRAIN windows at B1/C20480/D1536 provide another
12,288 labels on Qwen3-8B (36 layers). Both custom backends retain its native
per-head query/key normalization before RoPE. All arms use the same FP16
conversion of the BF16 checkpoint, not native BF16 inference. The shared HF
prefill uses 1,024-token chunks; recurrent labels use the native tokenizer with
no artificial BOS/EOS or chat template. Own HF PPL is 7.47057. Table~\ref{tab:qwen}
shows close development fidelity; PageGauge's interval includes one, not
evidence of improved predictive quality. Top-1 agreement is 99.683\%.
All eight runs consume 48 newly aged INT8 history pages per layer.

\begin{table*}[t]
\caption{Qwen3-8B development quality and allocated served KV storage.
Eight windows, 12,288 labels; descriptive paired-window intervals.
These are not speed, generated-task, independent TEST, or non-inferiority results.}
\label{tab:qwen}\centering\small
% Generated development evidence; do not edit by hand.
\begin{tabular}{lrrrrrr}
\toprule
Method & PPL & HF PPL & PPL/HF & Descriptive 95\% interval & KV MiB & KV reduction\\\midrule
FlashInfer FP16 & 7.46953 & 7.47057 & 0.99986 & [0.99972, 0.99995] & 3096.00 & 0.00\% \\
PageGauge INT8 & 7.46912 & 7.47057 & 0.99981 & [0.99945, 1.00012] & 1954.93 & 36.86\% \\
Kitty-Pro$^\dagger$ & 7.59687 & 7.46873 & 1.01716 & [1.00248, 1.03602] & 521.39 & 83.16\% \\
\bottomrule
\end{tabular}

\par\smallskip\footnotesize
$^\dagger$Native Kitty-Pro: two-bit K/V with 25\% of K channels boosted to
four bits. Two loaded channel-index operands are widened to int32 before
byte-offset arithmetic; stored metadata and quantization policy are unchanged.
Kitty uses its pinned Transformers 4.53.2 fork and full-prefix SDPA prefill;
PageGauge/FlashInfer use 4.57.6 and chunked prefill. Each row retains its own
HF reference on exactly matched token IDs. MiB denotes $2^{20}$ bytes.
\end{table*}

\begin{table*}[t]
\caption{CPU-backed common-engine timing pilot, B4/C20480/D1536. Step latency
is for the whole batch; throughput counts four output tokens per step.
Ranges are the minimum and maximum of three within-process repeats, not
confidence intervals. Peak allocation includes weights and native temporary
allocations during timed decode; GiB denotes $2^{30}$ bytes.}
\label{tab:frontier}\centering\small
% Generated development evidence; do not edit by hand.
\begin{tabular}{lrrrrr}
\toprule
Method & Step ms & Repeat range ms & Aggregate tokens/s & KV GiB & Peak alloc. GiB\\\midrule
FlashInfer FP16 & 22.11 & [20.78, 22.67] & 180.95 & 10.75 & 24.92 \\
PageGauge INT8 (split32) & 23.73 & [21.71, 24.94] & 168.56 & 6.79 & 21.09 \\
KIVI INT4 & 70.89 & [70.85, 70.89] & 56.43 & 3.36 & 21.37 \\
KIVI INT2 & 61.09 & [60.95, 62.18] & 65.48 & 2.02 & 18.48 \\
BitDecoding INT4 & 28.03 & [27.40, 31.36] & 142.69 & 3.42 & 17.58 \\
\bottomrule
\end{tabular}

\end{table*}

The matched native Kitty-Pro cohort reduces served KV storage by 83.16\%,
with PPL/HF 1.01716 and descriptive interval [1.00248, 1.03602]. Its top-1
agreement is 93.156\%. Thus it offers substantially smaller storage but greater
predictive distortion on these windows; this is a quality--memory trade-off,
not an equal-bit or equal-memory comparison. No Kitty speed claim follows.

Before this cohort, the native port failed the same-cache execution oracle.
The loaded uint8 boosted-channel indices overflowed when multiplied by a
32-byte address stride in the inspected compiled kernel. Widening both the
packing and attention address operands to int32 restored the intended layout.
The corrected native port passes B1/B4 recurrent same-cache checks, with maximum
relative $L_2$ error below $4.0\times10^{-5}$. This compatibility correction is
disclosed separately from predictive quality; the failed original-port
diagnostic is retained. It does not change PageGauge's method.

The preceding short C4096/D785 pilot also exposes a limitation: PageGauge
allocates 786,440,448 KV bytes versus 721,944,576 for FP16. Its fixed exact-region
storage and allocated code capacity outweigh compression at this context.
The long-context memory reduction therefore must not be generalized to every
sequence length. That pilot's PPL/HF ratio is 1.00006 on only 785 labels.

\section{Common-Engine Timing and Memory}
A separate Mistral experiment puts all five attention backends in the same
eager Transformers 4.36.2 body, without packed projection replacement or CUDA
graph capture. Four independent TRAIN requests use 20,480 prefill tokens and
1,536 decode steps. Each method gets one fresh process, one full recurrence
warmup, and three timed recurrences. This is a pilot, not a hierarchical
confidence-interval experiment or each method's best serving implementation.

Initial serving states are packed from serial coherent HF prefills, stored on
the CPU, and uploaded before each recurrence. No original GPU FP16 reference
cache is retained during decode. Prefill, packing, host storage, and upload are
outside the decode timing and are recorded separately. The common 28 GiB
PyTorch allocator budget does not cap allocations outside that allocator.
Served KV excludes BitDecoding's separately recorded append staging; peak
allocation includes it. These are not maximum-batch or admission-cost results.

PageGauge's median latency is higher than FlashInfer's in this common eager
body, with overlapping and variable repeat ranges (Table~\ref{tab:frontier}).
This negative point result is retained, not pooled with optimized-decoder
measurements. More attention launches and integration overhead are candidate
explanations to test, not established causes. A faster compressed-history kernel
alone would not establish an application-level speedup.

Native KIVI INT4 also illustrates why cache bytes alone do not determine
deployment capacity: its timed peak allocation is 22,941,229,568 bytes versus
18,684,578,304 bytes at the final boundary. Native GQA repeats, concatenation,
and old/new legacy-cache lifetimes are possible contributors, requiring
separate attribution. No native-baseline policy is changed to improve its
position in this completed cohort.

\clearpage
\section{Cross-Domain Development: PG-19}
Eight PG-19 training books are selected by a fixed hash ordering of object
names among metadata-size-eligible books, before contents or outcomes are
examined. One fixed B1/C20480/D1536 window per book supplies 12,288 labels per
model for Mistral-7B-v0.3 and Qwen3-8B. Books are shared across models, but their
native tokenizers yield different token sequences. The S4/A128/T768
representation is unchanged; PageGauge uses the exact32 launch selection.
Each trajectory consumes 48 newly aged INT8 history pages. There is no
quality-driven stopping or book replacement.

\begin{table}[ht]
\caption{Eight PG-19 TRAIN books per model. M7: Mistral-7B-v0.3; Q8: Qwen3-8B.
Own HF model-token PPL is 6.42158 and 12.08988, respectively.
Intervals are descriptive paired-book bootstrap intervals, not final TEST
or prospective non-inferiority results.}
\label{tab:books}\centering\footnotesize
% Generated development evidence; do not edit by hand.
\begin{tabular}{lrr}
\toprule
Method & PPL & PPL/HF [95\% interval]\\\midrule
M7 FI FP16 & 6.42145 & 0.99998 [0.99993, 1.00003] \\
M7 PG INT8 & 6.42155 & 0.99999 [0.99993, 1.00007] \\
Q8 FI FP16 & 12.09083 & 1.00008 [1.00001, 1.00014] \\
Q8 PG INT8 & 12.09432 & 1.00037 [0.99991, 1.00101] \\
\bottomrule
\end{tabular}

\end{table}

PageGauge retains 99.805\% and 99.634\% top-1 agreement for Mistral and Qwen,
respectively, and reduces served KV storage by 36.86\% for both. Qwen's PPL
point estimate increases by 0.0367\%; its interval permits an increase of about
0.1014\%. Both PageGauge PPL-ratio intervals include one. This supports close
development fidelity, not superior quality or a prospective non-inferiority pass.
These are subword/model-token PPL measurements on selected windows, not the
official whole-corpus word-normalized PG-19 score. The training books are now
exposed development data, and may also have appeared in model pretraining.
Independent final evaluation remains outstanding.

\section{Attention-Dispatch Integration Pilot}
An opt-in follow-up reuses the existing attention-only CUDA graph capture
for both FlashInfer and PageGauge, leaving the common dense body, kernels,
cache policies, append and planning unchanged. All attention topologies are
prepared before decode timing; there is no eager fallback. At
B4/C20480/D1536, each backend completes 49,152 graph replays. All 352 selected
same-cache attention comparisons per backend are exactly equal to the original
execution, and the restored cache is bitwise unchanged by capture.

\begin{table}[ht]
\caption{Matched attention-only graph pilot. One fresh process per arm,
one full recurrence warmup and three timed repeats on the same TRAIN workload.
Ranges are not confidence intervals.}
\label{tab:graphs}\centering\small
% Generated development evidence; do not edit by hand.
\begin{tabular}{lrr}
\toprule
Method & Step ms & Repeat range ms\\\midrule
FI FP16 & 21.08 & [20.95, 21.72] \\
PG INT8 & 20.72 & [20.66, 21.63] \\
\bottomrule
\end{tabular}

\end{table}

The point ratio is only 1.01731$\times$ FlashInfer/PageGauge, with overlapping
repeat ranges (Table~\ref{tab:graphs}). This does not establish a reliable
speedup or the aspirational 1.2--1.5$\times$ target. The earlier eager negative
point is retained as a separate dispatch configuration. Across restored rounds,
graph preparation takes 0.074--0.159\,s for FlashInfer and 0.342--0.422\,s for
PageGauge, outside decode timing. Each first capture increases device-used
memory by 16\,MiB. Order-balanced replication and setup amortization remain
necessary before a final integration claim; no production default is changed.

\section{Capacity Grid and Integration Diagnosis}
The original common-eager configuration is tested on a fixed B4/B8/B16 grid
at C20480/D1536 and the same 28\,GiB PyTorch allocator budget. The completed
B4 pilot is reused. New feasible points require a full recurrence warmup and
one full timed recurrence; these are capacity pilots, not replicated throughput
estimates. Larger points are excluded without execution only when FP16 parameter
storage plus a conservative minimum-bit live-cache bound already exceeds the
budget. Metadata, exact-region premiums and temporary allocations are omitted
from that bound. It is not an observed CUDA OOM.

\begin{table}[ht]
\caption{Original fixed capacity grid. OK: complete recurrence; OOM: measured
CUDA allocation failure; Bound: conservative analytical exclusion. These are
verified grid points, not global maximum batch sizes.}
\label{tab:capacity}\centering\small
% Generated development evidence; do not edit by hand.
\begin{tabular}{lccc}
\toprule
Method & B4 & B8 & B16\\\midrule
FI FP16 & OK & Bound & Bound \\
PG INT8 & OK & OOM & Bound \\
KIVI INT4 & OK & OOM & OOM \\
KIVI INT2 & OK & OK & OOM \\
BitDecoding INT4 & OK & OK & OOM \\
\bottomrule
\end{tabular}

\end{table}

At B8, BitDecoding INT4 completes at 27.75\,ms per step (288.24 aggregate
tokens/s), while KIVI INT2 completes at 87.05\,ms (91.90 tokens/s). No B16 point
completes under the original configuration. Failed outcomes remain in the
record, including PageGauge's B8 setup failure and the native warmup failures.

An explicit-default replay isolates PageGauge's failing 128\,MiB workspace
allocation. The allocator has 253.54\,MiB of inactive split storage across 67
blocks, but its largest inactive block is only 4\,MiB. Live storage plus the
requested workspace fits the budget; allocator reserve prevents a fresh
allocation. Enabling expandable segments removes that OOM without changing
the budget, representation or kernel math. This exposes a second issue:
graph-padded planning cannot accommodate the fixed split at B8, even though
the common eager benchmark never replays a graph.

A separate opt-in adapter therefore uses non-graph planning with the declared
expandable allocator. Its complete B8/C20480/D1536 validation checks 12,288
labels and preserves the restored cache bit-for-bit. Thirty-six selected
same-cache comparisons reconstruct one KV head/request at a time for VRAM:
maximum relative $L_2$ error is 0.00098 and maximum absolute error is 0.00192,
below unchanged execution limits of 0.005 and 0.02. These cover 144 query rows,
not every head at every step. Development PPL is 5.20478 versus its own serial
HF reference 5.20417 (ratio 1.000116); top-1 agreement is 99.829\%.

The original failures are not relabeled as successes. A separate fresh-process
B8 pilot completes one full recurrence warmup and one timed recurrence at
25.68\,ms per step, or 311.50 aggregate tokens/s. Its state restore and bitwise
checks take 23.85\,s outside decode timing; this is diagnostic CPU-backed setup,
not native-prefill or admission latency. All four originally failed native
cells also receive explicit-default and expandable-allocator runs: KIVI INT4
at B8/B16, and KIVI INT2 and BitDecoding INT4 at B16. All eight runs retain
measured CUDA OOM under the same budget, with no silent fallback. This
diagnostic does not establish a global capacity maximum or a throughput
advantage; timing replication remains necessary. No production default is
changed. The record also retains the
first failed validation of the new checker: its GQA-expanded center was
incorrectly indexed, then corrected and rerun without relaxing tolerances.

\section{Supporting Shift-Invariant Error Analysis}
The exact factorization evaluates a reconstructed cache; it does not remove
quantization error. For one fixed query and head, on the same unmasked support,
write $x_j=\alpha q^\top k_j$ and
$\widehat x_j=x_j+\delta_j$, with
$\delta_j=\alpha q^\top(\widehat k_j-k_j)$.
Let $p=\operatorname{softmax}(x)$, $\widehat p=\operatorname{softmax}(\widehat x)$,
$\Delta=\max_j\delta_j-\min_j\delta_j$, and
$D_V=\max_{i,j}\|v_i-v_j\|$ for any norm. Then
\begin{equation}
\begin{split}
\|\widehat o-o\|\leq{}&\sum_j\widehat p_j\|\widehat v_j-v_j\|\\
&+D_V\tanh(\Delta/4).
\end{split}
\label{eq:rangebound}
\end{equation}

To see this, set $R_j=e^{\delta_j}$, $Z=\mathbb E_p R$,
$a=\min R_j$, and $b=\max R_j$. If $a=b$, the distributions coincide.
Otherwise, bounding $|R-Z|$ by its chord on $[a,b]$ gives
\begin{equation}
\begin{split}
\operatorname{TV}(p,\widehat p)
&\leq \frac{(b-Z)(Z-a)}{(b-a)Z}\\
&\leq \frac{\sqrt b-\sqrt a}{\sqrt b+\sqrt a}
=\tanh(\Delta/4).
\end{split}
\end{equation}
The second inequality is maximized at $Z=\sqrt{ab}$. The positive and negative
parts of $\widehat p-p$ each have mass TV; their normalized value averages lie
in a convex hull of diameter at most $D_V$. Applying the triangle inequality
to $\widehat o-o=\sum_j\widehat p_j(\widehat v_j-v_j)
+\sum_j(\widehat p_j-p_j)v_j$ proves (\ref{eq:rangebound}). The TV constant
is attained by two logits $x=(\Delta/2,0)$ with perturbation
$\delta=(-\Delta/2,\Delta/2)$.

A common key-error vector shifts all logits equally and has $\Delta=0$;
page-dependent errors generally do not. Exact rows have zero representation
error at those rows, not necessarily elsewhere. If queries also differ, use
the total logit error
$\delta_j=\alpha(q^\top e_j^K+e_q^\top k_j+e_q^\top e_j^K)$.
Finite-precision execution error is a separate contribution. These bounds are
supporting analysis, not a new-algorithm novelty claim, a full-network guarantee,
or a substitute for measured task quality.

For final model logits $z$ and perturbations $e$, a top-1 margin larger than
$\max e-\min e$ is sufficient for token agreement, but agreement alone need
not imply small distributional error. For any target token $y$, the NLL change
is $\log\mathbb E_{\operatorname{softmax}(z)}e^{e_j}-e_y$ and has magnitude at
most $\max e-\min e$. Neither result yields a universal cosine threshold or
a speed guarantee. Historical gates remain historical; final evaluation still
requires an independently frozen protocol.


\section{Native Prompt and Decode Costs}
Table~\ref{tab:nativecost} adds native NSNQuant/Mistral and Kitty/Qwen3
request-cost pilots, each with its own Transformers HF control. Each arm has
one fresh process, a full-request warmup and three repeated B1/C20480/D1536
requests. Native cache allocation, quantization and full-prefix SDPA plus the
last prompt LM head are included in prefill; model loading and static weight
rotation are separate. Decode includes all layers, recurrent updates and the
LM head, using teacher-forced inputs without sampling or server overhead.
No GPU FP16 reference cache is retained in the native quantized arms.

\begin{table*}[t]
\caption{Native request-cost pilots on one exposed TRAIN fixture per model.
Medians and within-process repeat ranges, not confidence intervals. Peak is
the maximum allocated storage across prefill/decode repeats, including model
weights; served KV and its reduction are separate. Kitty-Pro was rerun after
an external interruption, while the three completed arms were retained.
The time gap precludes an adjacent matched Kitty speed comparison.}
\label{tab:nativecost}\centering\small
% Generated development evidence; do not edit by hand.
\begin{tabular}{llrrrrr}
\toprule
Model & Method & Prefill s & Step ms [range] & KV MiB & KV reduction & Peak GiB\\\midrule
Mistral & HF 4.48.1 & 2.172 & 32.50 [32.05, 32.63] & 2752.00 & 0.00\% & 18.31 \\
Mistral & NSN INT2 & 2.216 & 31.17 [28.66, 31.40] & 385.17 & 86.00\% & 16.16 \\
Qwen3 & HF 4.53.2 & 2.326 & 44.21 [43.63, 44.98] & 3096.00 & 0.00\% & 20.14 \\
Qwen3 & Kitty-Pro & 2.830 & 49.77 [49.70, 50.01] & 521.39 & 83.16\% & 17.84 \\
\bottomrule
\end{tabular}

\end{table*}

\begin{table*}[t]
\caption{Fixed regional-policy ablation on eight PG19 TRAIN books. Ratios
near one indicate similar model-token PPL, not a prospective non-inferiority
test. The interval is a paired-book descriptive bootstrap relative to the
reference PageGauge policy; top-1 agreement is instead relative to HF.
KV reduction uses the same full FP16 cache as the denominator for every row.}
\label{tab:regional}\centering\small
% Generated development evidence; do not edit by hand.
\begin{tabular}{lrrrrr}
\toprule
Policy & PPL/HF & PPL/reference policy [95\% interval] & Top1/HF \% & KV MiB & KV reduction\\\midrule
S4/A128/T768 & 0.999995 & 1.000000 [1.000000, 1.000000] & 99.805 & 1737.72 & 36.86\% \\
S0/A128/T768 & 1.000336 & 1.000341 [0.999565, 1.001139] & 99.349 & 1729.72 & 37.15\% \\
S4/A0/T768 & 0.999910 & 0.999915 [0.999849, 0.999991] & 99.805 & 1481.72 & 46.16\% \\
S4/A128/T16 & 0.999941 & 0.999946 [0.999807, 1.000092] & 99.666 & 1643.72 & 40.27\% \\
\bottomrule
\end{tabular}

\end{table*}

Lower cache storage does not imply an equal reduction in request cost.
NSNQuant's fixed-B1 decoding gain over its HF point is modest, while its
prefill is slightly longer. Kitty-Pro's observed decode and prefill costs are
higher than the earlier HF point, but the interruption prevents attributing
that difference solely to its method. These results neither measure
memory-enabled maximum throughput nor establish a PageGauge-relative advantage
across different engines. Balanced fresh-process replication and actual
serving measurements remain necessary.

\section{Exact-Region Ablation}
We evaluate four fixed, one-factor policies on the same eight preselected
Mistral PG19 TRAIN books, with 1536 recurrent labels per book and policy.
All policies share the coherent HF prefill, INT8 quantizer, scalar metadata,
centers and kernels. The reference is S4/A128/T768; the alternatives remove
the exact prefix, remove the static prefill suffix, or reduce the recent tail
to one page. T16 is the minimum supported nonempty page, not a zero-residual
configuration. No production default is changed by this development ablation.

Selected same-cache FP32 attention checks cover three layers, three decode
positions and two KV-head groups per book and policy: 144 group checks per
policy. These check execution of the reconstructed cache, not fidelity to
unquantized HF activations. Attention-region mass is measured only for these
selected queries using each candidate's reconstructed keys, not all model
attention or the original HF distribution. Changing an exact region changes
the representation error even though the factorization identity is unchanged.
Table~\ref{tab:regional} reports the predictive and storage consequences.
Removing the static suffix reduces PageGauge storage by 14.73\% while
retaining the reference policy's 99.805\% top-1 agreement with HF on this
cohort. Its PPL/reference-policy ratio is 0.999915, a very small change;
this exploratory multi-policy comparison is not evidence of general quality
superiority. Removing the prefix instead lowers top-1 agreement to 99.349\%
for only 0.46\% less PageGauge storage. Cross-model confirmation is pending.
The instrumented quality loop is not used as a latency measurement; clean
fresh-process costs and replication remain necessary before policy selection.

\section{Outstanding Evidence}
Remaining work includes generated-task quality, additional model families,
shape/GPU coverage, residual-policy timing, replicated post-diagnosis
throughput, capacity boundaries, and realistic serving with prefill and
admission. Native cost pilots are not replicated serving evidence. The completed
capacity grid is not a global capacity search or an optimized serving result.
Fixed value conditioning is not promoted: its development results did not
establish a useful quality improvement. Final code and input policies must
be frozen before independent TEST exposure.

This is not the full main manuscript; submitted workshop artifacts are
unchanged. Tables use hash-verified raw results, with provenance in the
accompanying local evidence manifest.
\end{document}
